% =============================================================================
% style/base_chapter.tex
% =============================================================================


% =============================================================================
% Двузначный номер главы в бейдже; в приложениях оставляем A, B, ...
% =============================================================================

\makeatletter
\@ifundefined{ifchapterbadgeappendix}{\newif\ifchapterbadgeappendix}{}
\makeatother

\pretocmd{\appendix}{\chapterbadgeappendixtrue}{}{}

\providecommand{\chapterBadgeNumber}{}
\renewcommand{\chapterBadgeNumber}{%
  \ifchapterbadgeappendix%
    \thechapter%
  \else%
    \ifnum\value{chapter}<10 0\fi\arabic{chapter}%
  \fi%
}


% =============================================================================
% Цвета заголовков глав
% =============================================================================

% Нумерованный шестиугольный бейдж.
\colorlet{chapterBadgeTop}{xmatrix!55!white}
\colorlet{chapterBadgeBottom}{xmatrix!85!black}
\colorlet{chapterBadgeFrame}{xmatrix!45!black}
\colorlet{chapterBadgeText}{xmatrix!3!white}

% Одна горизонтальная линия под нумерованной главой.
\colorlet{chapterRuleColor}{xmatrix!60!black}

% Толстая вертикальная метка для ненумерованной главы.
\colorlet{chapterNumberlessLine}{xmatrix!60!black}


% =============================================================================
% Fading для чистого верхнего блика
% =============================================================================

\tikzfading[
  name=chapterGlossFade,
  top color=transparent!0,
  bottom color=transparent!100,
]


% =============================================================================
% Геометрия и выравнивание декоративных элементов заголовка
% =============================================================================

\newlength{\chapterBadgeBaselineShift}
\setlength{\chapterBadgeBaselineShift}{6.3pt}

\newlength{\chapterNumberlessBaselineShift}
\setlength{\chapterNumberlessBaselineShift}{6.3pt}

% Контур шестиугольного бейджа.
\newcommand{\chapterBadgeHexPath}[1]{%
  ([xshift=2.8mm]#1.north west) --
  ([xshift=-2.8mm]#1.north east) --
  ([yshift=-0.25mm]#1.east) --
  ([xshift=-2.8mm]#1.south east) --
  ([xshift=2.8mm]#1.south west) --
  ([yshift=0.25mm]#1.west) --
  cycle%
}


% =============================================================================
% Бейдж нумерованной главы
% =============================================================================

\providecommand{\chapterMatrixBadgeE}[1]{}
\renewcommand{\chapterMatrixBadgeE}[1]{%
  \tikz[baseline={([yshift=-\chapterBadgeBaselineShift]B.center)}]{%
    % Невидимый узел задаёт размеры бейджа.
    \node[
      minimum width=18.4mm,
      minimum height=12.0mm,
      inner xsep=5.8pt,
      inner ysep=1.9pt,
      outer sep=0pt,
      font=\sffamily\bfseries\fontsize{20}{22}\selectfont
    ] (B) {\phantom{\strut #1}};

    % Внешняя мягкая тень.
    \begin{scope}[xshift=0.50mm, yshift=-0.60mm]
      \path[
        fill=black,
        opacity=0.18,
        rounded corners=1.8pt,
      ]
        \chapterBadgeHexPath{B};
    \end{scope}

    % Основная заливка: чистый вертикальный градиент без среднего цвета.
    \shade[
      top color=chapterBadgeTop,
      bottom color=chapterBadgeBottom,
      draw=chapterBadgeFrame,
      line width=0.80pt,
      rounded corners=1.8pt,
      line join=round,
    ]
      \chapterBadgeHexPath{B};

    % Верхний глянцевый блик: плавно исчезает к середине бейджа.
    \begin{scope}
      \clip[rounded corners=1.8pt]
        \chapterBadgeHexPath{B};
      \path[
        fill=white,
        path fading=chapterGlossFade,
        opacity=0.55,
      ]
        ([yshift=-0.25mm]B.north west) rectangle
        ([yshift=-0.25mm]B.east);
    \end{scope}

    % Тонкий внутренний световой кант вверху.
    \begin{scope}
      \clip[rounded corners=1.8pt]
        \chapterBadgeHexPath{B};
      \draw[
        white,
        opacity=0.55,
        line width=0.50pt,
        line cap=round,
      ]
        ([xshift=3.0mm, yshift=-0.42mm]B.north west) --
        ([xshift=-3.0mm, yshift=-0.42mm]B.north east);
    \end{scope}

    % Тонкий внутренний теневой кант внизу.
    \begin{scope}
      \clip[rounded corners=1.8pt]
        \chapterBadgeHexPath{B};
      \draw[
        black,
        opacity=0.22,
        line width=0.45pt,
        line cap=round,
      ]
        ([xshift=3.0mm, yshift=0.45mm]B.south west) --
        ([xshift=-3.0mm, yshift=0.45mm]B.south east);
    \end{scope}

    % Наружная обводка поверх всех эффектов.
    \path[
      draw=chapterBadgeFrame,
      line width=0.80pt,
      rounded corners=1.8pt,
      line join=round,
    ]
      \chapterBadgeHexPath{B};

    % Номер главы.
    \node[
      minimum width=18.4mm,
      minimum height=12.0mm,
      inner xsep=5.8pt,
      inner ysep=1.9pt,
      outer sep=0pt,
      text=chapterBadgeText,
      font=\sffamily\bfseries\fontsize{20}{22}\selectfont,
    ] at (B.center) {\strut #1};
  }%
}


% =============================================================================
% Метка ненумерованной главы: одна толстая вертикальная линия
% =============================================================================

\providecommand{\chapterNumberlessMark}{}
\renewcommand{\chapterNumberlessMark}{%
  \tikz[baseline={([yshift=-\chapterNumberlessBaselineShift]C.center)}]{%
    \node[
      minimum width=5.0mm,
      minimum height=12.0mm,
      inner sep=0pt,
      outer sep=0pt
    ] (C) {};

    \draw[
      chapterNumberlessLine,
      line width=2.6pt,
      line cap=round,
    ]
      ([xshift=2.2mm, yshift=-0.3mm]C.north west) --
      ([xshift=2.2mm, yshift=0.3mm]C.south west);
  }%
}


% =============================================================================
% Нумерованные главы: бейдж + одна толстая горизонтальная линия
% =============================================================================

\titleformat{\chapter}[hang]
  {%
    \normalfont\sffamily\bfseries\raggedright\color{xmatrixDark}%
    \hyphenpenalty=10000\exhyphenpenalty=10000%
  }
  {\chapterMatrixBadgeE{\chapterBadgeNumber}}
  {13pt}
  {\fontsize{24}{28}\selectfont}
  [\vspace{7.5pt}{\color{chapterRuleColor}\titlerule[1.90pt]}]


% =============================================================================
% Ненумерованные главы: только вертикальная метка, без горизонтальной линии
% =============================================================================

\titleformat{name=\chapter,numberless}[hang]
  {%
    \normalfont\sffamily\bfseries\raggedright\color{xmatrixDark}%
    \hyphenpenalty=10000\exhyphenpenalty=10000%
  }
  {\chapterNumberlessMark}
  {13pt}
  {\fontsize{24}{28}\selectfont}


% =============================================================================
% Отступы заголовка главы
% =============================================================================

\titlespacing*{\chapter}{0pt}{10pt plus 2pt minus 2pt}{20pt}